\documentclass[10pt,a4paper]{amsart}
\usepackage{amsmath,amssymb,amsthm,graphicx}
% Retained preamble for the Stacks Project chapter {CH} (shared project preamble).
% The chapter file pinned-original-source.tex only carries \input{{preamble}} on its first line,
% so no preamble text lives inside the chapter itself.  The Stacks Project uses one single shared
% preamble for every chapter: the snapshot stages it three times (sources/stacks-chow/preamble.tex,
% sources/stacks-etale-cohomology/preamble.tex and focused-seeds/AGP/sources/stacks-snapshot-20261011/
% preamble.tex) and those copies are identical in size.  The block below is copied mechanically,
% line for line, from sources/stacks-chow/preamble.tex lines 14-18 and 154-217: the Xy-pic setup and
% the native theorem declarations together with the mathematical macros.  Nothing was added, removed
% or redefined; the wrapper class and the amsmath/amssymb/amsthm/graphicx base packages are supplied
% by the assembler.
% For commutative diagrams we use Xy-pic
% Chapter: stacks-spaces-more-morphisms.tex -- More on Morphisms of Spaces.
\usepackage[all]{xy}

% We use 2cell for 2-commutative diagrams.
\xyoption{2cell}

% Theorem environments.
%
\theoremstyle{plain}
\newtheorem{theorem}[subsection]{Theorem}
\newtheorem{proposition}[subsection]{Proposition}
\newtheorem{lemma}[subsection]{Lemma}

\theoremstyle{definition}
\newtheorem{definition}[subsection]{Definition}
\newtheorem{example}[subsection]{Example}
\newtheorem{exercise}[subsection]{Exercise}
\newtheorem{situation}[subsection]{Situation}

\theoremstyle{remark}
\newtheorem{remark}[subsection]{Remark}
\newtheorem{remarks}[subsection]{Remarks}

\numberwithin{equation}{subsection}

% Macros
%
\def\lim{\mathop{\mathrm{lim}}\nolimits}
\def\colim{\mathop{\mathrm{colim}}\nolimits}
\def\Spec{\mathop{\mathrm{Spec}}}
\def\Hom{\mathop{\mathrm{Hom}}\nolimits}
\def\Ext{\mathop{\mathrm{Ext}}\nolimits}
\def\SheafHom{\mathop{\mathcal{H}\!\mathit{om}}\nolimits}
\def\SheafExt{\mathop{\mathcal{E}\!\mathit{xt}}\nolimits}
\def\Sch{\mathit{Sch}}
\def\Mor{\mathop{\mathrm{Mor}}\nolimits}
\def\Ob{\mathop{\mathrm{Ob}}\nolimits}
\def\Sh{\mathop{\mathit{Sh}}\nolimits}
\def\NL{\mathop{N\!L}\nolimits}
\def\CH{\mathop{\mathrm{CH}}\nolimits}
\def\proetale{{pro\text{-}\acute{e}tale}}
\def\etale{{\acute{e}tale}}
\def\QCoh{\mathit{QCoh}}
\def\Ker{\mathop{\mathrm{Ker}}}
\def\Im{\mathop{\mathrm{Im}}}
\def\Coker{\mathop{\mathrm{Coker}}}
\def\Coim{\mathop{\mathrm{Coim}}}

% Boxtimes
%
\DeclareMathSymbol{\boxtimes}{\mathbin}{AMSa}{"02}

%
% Macros for moduli stacks/spaces
%
\def\QCohstack{\mathcal{QC}\!\mathit{oh}}
\def\Cohstack{\mathcal{C}\!\mathit{oh}}
\def\Spacesstack{\mathcal{S}\!\mathit{paces}}
\def\Quotfunctor{\mathrm{Quot}}
\def\Hilbfunctor{\mathrm{Hilb}}
\def\Curvesstack{\mathcal{C}\!\mathit{urves}}
\def\Polarizedstack{\mathcal{P}\!\mathit{olarized}}
\def\Complexesstack{\mathcal{C}\!\mathit{omplexes}}
% \Pic is the operator that assigns to X its picard group, usage \Pic(X)
% \Picardstack_{X/B} denotes the Picard stack of X over B
% \Picardfunctor_{X/B} denotes the Picard functor of X over B
\def\Pic{\mathop{\mathrm{Pic}}\nolimits}
\def\Picardstack{\mathcal{P}\!\mathit{ic}}
\def\Picardfunctor{\mathrm{Pic}}
\def\Deformationcategory{\mathcal{D}\!\mathit{ef}}
\title{Formal unramified, etale and smooth properties cancel on the left of a formally unramified transformation}
\author{The Stacks Project Authors}
\date{stacks-project:spaces-more-morphisms.tex:local-snapshot-20261011; GNU FDL 1.2 or later}
\begin{document}
\maketitle
\noindent\textit{Editorial note.} Statement span is the whole lemma environment lines 2665-2676, including the three enumerated cases (1) unramified, (2) etale, (3) smooth.\par\medskip
\noindent\textit{Editorial note.} Proof span is the complete author proof, lines 2678-2691; the closing sentence 'From this the lemma follows formally.' is retained verbatim and is the author's own reduction, not a gap in the retained argument.\par\medskip
\noindent\textit{Editorial note.} No external\_references are needed: the proof cites no label and no bibliography entry, and the statement cites nothing.\par\medskip
\noindent\textit{Editorial note.} Preamble is the retained chapter preamble (sources/stacks-spaces-more-morphisms/retained-preamble.tex) copied mechanically from the native project preamble; it stops before the document body and contains no documentclass, so the assembler supplies the wrapper class.\par\medskip
\noindent\textit{Editorial note.} extra\_packages is empty on purpose: the retained preamble already loads \textbackslash{}usepackage[all]\{xy\}, and declaring xy again in extra\_packages produces an option clash in the assembler wrapper. This matches the 303 sibling cards that reuse a retained or native preamble.\par\medskip
\noindent\textit{Editorial note.} No figures or bibliography assets are referenced by this span.\par\medskip
\noindent\textit{Allowed standard prerequisites (editorial).}\begin{itemize}
\item Explicit prerequisite boundary. Definitions of formally unramified, formally etale and formally smooth transformations of functors on (Sch/S)\_fppf, and the square-zero thickening lifting criterion they are defined by. The correspondence argument in the retained author proof is not replaced by any cited result.
\end{itemize}\medskip
\begin{lemma}
\label{lemma-formally-permanence}
Let $S$ be a scheme.
Let $F, G, H : (\Sch/S)_{fppf}^{opp} \to \textit{Sets}$.
Let $a : F \to G$, $b : G \to H$ be transformations of functors.
Assume $b$ is formally unramified.
\begin{enumerate}
\item If $b \circ a$ is formally unramified then $a$ is formally unramified.
\item If $b \circ a$ is formally \'etale then $a$ is formally \'etale.
\item If $b \circ a$ is formally smooth then $a$ is formally smooth.
\end{enumerate}
\end{lemma}

\begin{proof}
Let $T \subset T'$ be a closed immersion of affine schemes defined by an ideal
of square zero. Let $g' : T' \to G$ and $f : T \to F$ be given such that
$g'|_T = a \circ f$. Because $b$ is formally unramified, there is a one
to one correspondence between
$$
\{f' : T' \to F \mid f = f'|_T\text{ and }a \circ f' = g'\}
$$
and
$$
\{f' : T' \to F \mid f = f'|_T\text{ and }b \circ a \circ f' = b \circ g'\}.
$$
From this the lemma follows formally.
\end{proof}

\end{document}
